%================================================
% 数学分析习题集题目环境
%
% 不使用 tcolorbox 绘制题目主体，
% 仅用 tcolorbox 绘制元信息框。
%================================================

\newcounter{problem}[chapter]

\renewcommand{\theproblem}{%
    \thechapter.\arabic{problem}%
}

%================================================
% 掌握程度文字
%================================================

\expandafter\def
\csname mastery@text@mastered\endcsname
{完全会}

\expandafter\def
\csname mastery@text@familiar\endcsname
{熟悉}

\expandafter\def
\csname mastery@text@unfamiliar\endcsname
{不熟悉}

\expandafter\def
\csname mastery@text@unknown\endcsname
{完全不会}

\expandafter\def
\csname mastery@text@unrated\endcsname
{未评定}

%================================================
% 掌握程度颜色
%================================================

\expandafter\def
\csname mastery@color@mastered\endcsname
{MasteredColor}

\expandafter\def
\csname mastery@color@familiar\endcsname
{FamiliarColor}

\expandafter\def
\csname mastery@color@unfamiliar\endcsname
{UnfamiliarColor}

\expandafter\def
\csname mastery@color@unknown\endcsname
{UnknownColor}

\expandafter\def
\csname mastery@color@unrated\endcsname
{UnratedColor}

\newcommand{\MasteryText}[1]{%
    \ifcsname mastery@text@#1\endcsname
        \csname mastery@text@#1\endcsname
    \else
        未评定%
    \fi
}

\newcommand{\MasteryColor}[1]{%
    \ifcsname mastery@color@#1\endcsname
        \csname mastery@color@#1\endcsname
    \else
        UnratedColor%
    \fi
}

\newcommand{\CurrentProblemNumber}{}
\newcommand{\CurrentProblemCode}{}

%================================================
% 元信息框（tcolorbox）
%
% 极浅底色 + 细边框，内容一行横排。
%================================================

\tcbuselibrary{skins}

\newcommand{\probleminfobox}[4]{%
    % #1 掌握程度 key，#2 编号，#3 难度，#4 题源
    \begin{tcolorbox}[
        enhanced,
        boxrule=0.5pt,
        colframe=grayblue!40,
        colback=PaleSky,
        sharp corners,
        left=6pt,
        right=6pt,
        top=3pt,
        bottom=3pt,
        boxsep=0pt,
    ]
    \small
    % 彩色小圆点 + 掌握程度文字
    \textcolor{\MasteryColor{#1}}{$\bullet$}%
    \,\textcolor{\MasteryColor{#1}}{\MasteryText{#1}}%
    %
    \if\relax\detokenize{#3}\relax\else
        \quad
        \textcolor{grayblue}{难度：#3}%
    \fi
    %
    \if\relax\detokenize{#2}\relax\else
        \quad
        \textcolor{grayblue}{编号：#2}%
    \fi
    %
    \if\relax\detokenize{#4}\relax\else
        \quad
        \textcolor{grayblue}{题源：#4}%
    \fi
    \end{tcolorbox}%
}

%================================================
% 题目环境
%
% 用法：
%
% \begin{problem}
%   [mastered]
%   {MA-P000001}
%   {题目名称}
%   {难度}
%   {题源}
%
% 题目正文
%
% \end{problem}
%================================================

\newenvironment{problem}[5][unrated]{%

    \refstepcounter{problem}%

    \xdef\CurrentProblemNumber{\theproblem}%
    \gdef\CurrentProblemCode{#2}%

    \par
    \addvspace{1.25em}%

    % 题号 + 题名标题行
    \noindent
    {%
        \large\bfseries
        题~\theproblem\quad #3%
        \par
    }%

    \vspace{0.3em}%

    % 元信息框
    \probleminfobox{#1}{#2}{#4}{#5}%

    \vspace{0.45em}%

}{%

    \par
    \addvspace{0.35em}%

}

%================================================
% 解答环境
%================================================

\newenvironment{problemsolution}{%

    \par
    \addvspace{0.35em}%

    \proofFont
    \noindent
    \textbf{解.}\quad

}{%

    \hfill$\square$%

    \par
    \addvspace{0.65em}%

}

%================================================
% 注记环境
%================================================

\newenvironment{problemnote}{%

    \par
    \addvspace{0.35em}%

    \small
    \noindent
    \textbf{注.}\quad

}{%

    \par
    \normalsize
    \addvspace{0.35em}%

}

%================================================
% 掌握程度图例
%================================================

\newcommand{\masterylegend}{%
    \noindent
    \textcolor{MasteredColor}{$\bullet$}~\textcolor{MasteredColor}{完全会}%
    \qquad
    \textcolor{FamiliarColor}{$\bullet$}~\textcolor{FamiliarColor}{熟悉}%
    \qquad
    \textcolor{UnfamiliarColor}{$\bullet$}~\textcolor{UnfamiliarColor}{不熟悉}%
    \qquad
    \textcolor{UnknownColor}{$\bullet$}~\textcolor{UnknownColor}{完全不会}%
    \qquad
    \textcolor{UnratedColor}{$\bullet$}~\textcolor{UnratedColor}{未评定}%
}